%! AP Statistics — Unit 8, Lesson 8.3 — Guided Notes
\documentclass[10pt]{article}
\usepackage{apstats-article}
\usepackage{apstats-boxes}

\begin{document}

\pageheader{Unit 8, Lesson 8.3}{Guided Notes}
\namedateperiod

\begin{objectivebox}
By the end of this lesson, I will be able to:
\begin{itemize}
    \item[$\square$] \writeline
    \item[$\square$] \writeline
    \item[$\square$] \writeline
\end{itemize}
\end{objectivebox}

\begin{vocabbox}
\termblanklong{Chi-square test statistic ($\chi^2$)}
\termblanklong{Degrees of freedom (GOF)}
\termblanklong{P-value (chi-square GOF)}
\termblanklong{Conclusion template}
\end{vocabbox}

\begin{notesbox}{The Chi-Square Test Statistic}
When carrying out a chi-square goodness-of-fit test, we measure how far the observed
counts are from the expected counts using the \textbf{chi-square test statistic}:

\vspace{0.1in}
\[
    \chi^2 = \sum \frac{(\text{\blank{2.5cm}} - \text{\blank{2.5cm}})^2}{\text{\blank{2.5cm}}}
\]

\vspace{0.05in}
Each term $\dfrac{(O-E)^2}{E}$ measures \blank{8cm}.

\vspace{0.1in}
A \textbf{larger} $\chi^2$ value provides \blank{4cm} evidence against $H_0$.

\vspace{0.1in}
\textbf{Degrees of freedom} for a GOF test: \quad $df = $ \blank{2cm} $-$ \blank{2cm}

where the first blank is the \blank{4cm} of categories.
\end{notesbox}

\begin{notesbox}{Worked Example: Flower Color Distribution}
A botanist claims a flower comes in 4 colors in a 3:2:2:1 ratio. A random sample
of 160 flowers is collected.

\renewcommand{\arraystretch}{2.2}
\begin{center}
\begin{tabular}{lcccc}
    \textbf{Color}           & Red    & Pink   & White  & Yellow \\
    \hline
    \textbf{Null proportion} & 0.375  & 0.25   & 0.25   & 0.125  \\
    \textbf{Expected ($E$)}  & \blank{1.8cm} & \blank{1.8cm} & \blank{1.8cm} & \blank{1.8cm} \\
    \textbf{Observed ($O$)}  & 65     & 44     & 38     & 13     \\
    $\boldsymbol{(O-E)^2/E}$ & \blank{1.8cm} & \blank{1.8cm} & \blank{1.8cm} & \blank{1.8cm} \\
\end{tabular}
\end{center}

\vspace{0.1in}
$\chi^2 = $ \blank{1.5cm} $+$ \blank{1.5cm} $+$ \blank{1.5cm} $+$ \blank{1.5cm}
        $=$ \blank{2cm}

\vspace{0.1in}
$df = $ \blank{2cm} $- 1 =$ \blank{1.5cm}
\end{notesbox}

\begin{notesbox}{Finding and Interpreting the P-Value}
Using a chi-square table with $df =$ \blank{1.5cm}, look up $\chi^2 \approx$ \blank{2cm}.

\vspace{0.1in}
The p-value is \blank{2cm} (always a \textbf{right-tailed} area).

\vspace{0.15in}
\textbf{Interpret the p-value:}\\[4pt]
``The probability, given \blank{7cm} is true, of obtaining
a chi-square statistic of \blank{2cm} or \blank{2cm} is approximately \blank{2cm}.''

\vspace{0.1in}
\textit{Note:} The p-value is \emph{not} the probability that $H_0$ is true.
\end{notesbox}

\begin{notesbox}{Stating a Complete Conclusion}
Compare the p-value to the significance level $\alpha = 0.05$:

\vspace{0.1in}
``Because $p =$ \blank{2cm} \blank{1cm} $\alpha = 0.05$, we \blank{3.5cm} $H_0$.
We \blank{3cm} convincing evidence that \blank{7cm}.''

\vspace{0.1in}
\textbf{For the botanist example:} \writelines{2}
\end{notesbox}

\begin{practicebox}
\textbf{Quick Check:} Suppose a different sample of 160 flowers gave $\chi^2 = 9.20$
with $df = 3$. Using the chi-square table, is $p < 0.05$ or $p > 0.05$?
What conclusion would you reach at $\alpha = 0.05$?

\writelines{3}
\end{practicebox}

\end{document}
